\subsection{Refactorización de la Arquitectura Interna}\label{subsec:resultado-b}

El segundo objetivo buscaba refactorizar la arquitectura interna del \textit{frontend} y del \textit{backend} mediante patrones modulares que aislaran responsabilidades y habilitaran la evolución independiente de cada dominio funcional. La evidencia del resultado es de carácter estructural y se aprecia en la organización final del código.

En el \textit{frontend}, la aplicación quedó organizada en nueve \textit{bounded contexts} implementados como Nuxt Layers, con el tema visual aislado en una capa independiente, tal como se describió en la \autoref{subsec:arquitectura}. La \autoref{fig:resultado-feature-tree} muestra la estructura resultante de \textit{features}, donde cada dominio constituye una unidad autónoma y componible que puede modificarse, probarse o reemplazarse de forma independiente.

\newcommand\resultadoFeatureTreeCaption{Estructura final del \textit{frontend} organizada en \textit{bounded contexts} mediante Nuxt Layers. \hspace{1em}}
\begin{figure}[H]
  \centering
  \begin{verbatim}
paralegales-ui/src
|-- app-core
|   |-- layouts
|   |-- middleware
|   |-- navigation
|   |-- plugins
|   `-- public
|-- common
|   |-- assets
|   |-- components
|   |-- composables
|   |-- models
|   |-- server
|   |-- stores
|   `-- utils
|-- features
|   |-- active-tracking
|   |-- auth
|   |-- home
|   |-- my-processes
|   |-- payments
|   |-- process-details
|   |-- search-processes
|   |-- terms-and-conditions
|   `-- user-settings
`-- vendor
    `-- vuexy
  \end{verbatim}
  \caption[\resultadoFeatureTreeCaption]{\resultadoFeatureTreeCaption Fuente: Elaboración propia.}
  \label{fig:resultado-feature-tree}
\end{figure}

En el \textit{backend}, las funciones Lambda adoptaron la separación \textit{Handler--Service--Repository} con repositorios definidos por interfaz, expuestas de forma unificada a través de Amazon API Gateway (\autoref{fig:api-gateway}). Este desacoplamiento habilita tanto la sustitución del motor de persistencia sin tocar la lógica de negocio como la validación de dicha lógica de forma aislada; la verificación de ese desacoplamiento mediante pruebas unitarias con implementaciones \textit{mock} del repositorio se documenta en el capítulo dedicado a las pruebas del sistema.
